%% un-sdg.tex: UN Sustainable Development Goals (FEUP regulations)
%% -------------------------------------------------------

\chapter*{UN Sustainable Development Goals} \label{chap:UnitedNations}

The United Nations Sustainable Development Goals (SDGs) provide a global framework to achieve a better and more sustainable future for all. 
It includes 17 goals\footnote{\url{https://sdgs.un.org/goals}} to address the world's most pressing challenges, including poverty, inequality, climate change, environmental degradation, peace, and justice. 

The specific Sustainable Development Goals mentioned have the following names:

\begin{description}
\item [SDG 7]
Ensure access to affordable, reliable, sustainable and modern energy for all
\item [SDG 9]
Build resilient infrastructure, promote inclusive and sustainable industrialization and foster innovation
\end{description}

\begin{center}
\begin{tabular}{|l|l|p{58mm}|p{52mm}|}
\hline
\textbf{SDG} & \textbf{Target} & \textbf{Contribution} & \textbf{Performance Indicators and Metrics} \\ 
\hline
\hline
    \multirow{3}{*}{7} & 7.3 & The characterization and consideration of transmitter power consumption contribute to the development of energy-conscious high-speed communication electronics. Maintaining multi-gigabit data transmission while controlling circuit power is relevant to improving the energy efficiency of electronic communication systems. & Post-layout power consumption of 24.42 mW for EQ3 at the nominal condition; 10 Gb/s operating data rate; comparison of power across operating configurations/corners. \\
\hline
    \multirow{2}{*}{9} & 9.5 & The development and physical implementation of a configurable 10 Gb/s CMOS transmitter contributes to research and technological innovation in high-speed integrated circuits and communication interfaces. The work investigates programmable equalization and termination techniques and validates their operation after parasitic extraction and across process and operating variations. & 10 Gb/s data rate; complete schematic-to-layout implementation; USB 3.2 Gen 2 electrical-performance evaluation; operation across PVT and resistor corners; \(176\,\mu\mathrm{m}\times82\,\mu\mathrm{m}\) layout area. \\
\hline
\end{tabular}
\end{center}

This work contributes primarily to SDG 9 -- Industry, Innovation and Infrastructure, Target 9.5, by advancing research and technological development in high-speed integrated communication circuits. The design, physical implementation, and post-layout verification of a configurable 10 Gb/s CMOS transmitter demonstrate the application of integrated-circuit design techniques to improve the performance and robustness of modern serial interfaces. Its contribution can be assessed through technical metrics such as the achieved 10 Gb/s data rate, the complete schematic-to-layout implementation, the 176 µm $\times$ 82 µm occupied area, and the verification of equalization, termination, and signal integrity across process and operating variations. The work also relates to SDG 7 -- Affordable and Clean Energy, Target 7.3, through its consideration of energy efficiency in high-speed electronic communication. Although the transmitter is not directly associated with renewable-energy generation, reducing the energy required for high-data-rate communication is relevant to the broader objective of improving the efficiency of electronic systems. In this context, the principal performance metric is the transmitter power consumption, with the representative post-layout EQ3 configuration consuming 24.42 mW at 10 Gb/s, corresponding to approximately 2.44 pJ/bit. Together, these metrics provide measurable indicators of the work's contribution to technological innovation while also considering the energy cost associated with increasing communication bandwidth.